\SourceLeaf{224}
\fancyhead[L]{\small ALLGEMEINE UNTERSUCHUNGEN ÜBER GEKRÜMMTE FLÄCHEN}
\renewcommand{\Article}[1]{\par\medskip\Needspace{4\baselineskip}\pdfbookmark[1]{Art. #1}{superficies-art-#1}\centerline{\textbf{#1.}}\nobreak\smallskip}
\begin{center}
ALLGEMEINE UNTERSUCHUNGEN\par\medskip
{\large ÜBER GEKRÜMMTE FLÄCHEN.}\par\medskip
\rule{.65in}{.3pt}
\end{center}
\Article{1}
\indent Untersuchungen über die Richtungen verschiedener Geraden im Raum gewinnen meist erheblich an Klarheit und Einfachheit, wenn man eine Kugeloberfläche mit dem Radius \(=1\) um einen beliebigen Mittelpunkt zu Hilfe nimmt. Jeder ihrer Punkte soll die Richtung der Geraden darstellen, die zu dem in diesem Punkt endenden Radius parallel sind. Da die Lage jedes Punktes im Raum durch drei Koordinaten bestimmt ist, nämlich durch seine Abstände von drei festen, paarweise aufeinander senkrecht stehenden Ebenen, sind zunächst die Richtungen der zu diesen Ebenen senkrechten Achsen zu betrachten. Die Punkte der Kugeloberfläche, die diese Richtungen darstellen, bezeichnen wir mit \((1),(2),(3)\). Der Abstand zwischen je zwei dieser Punkte beträgt also einen Quadranten. Für die Achsen setzen wir außerdem die Richtungen voraus, in denen die entsprechenden Koordinaten wachsen.

\Article{2}
\indent Es dürfte nützlich sein, hier einige Sätze zusammenzustellen, die bei solchen Fragen häufig verwendet werden.

I. Der Winkel zwischen zwei sich schneidenden Geraden wird durch den Bogen zwischen den Punkten der Kugeloberfläche gemessen, die den Richtungen der Geraden entsprechen.

II. Die Lage einer beliebigen Ebene lässt sich durch einen Großkreis auf der Kugeloberfläche darstellen, dessen Ebene zur gegebenen Ebene parallel ist.
